% !TeX root = ../main.tex

\section{Вывод ядер Вольтерра}
\label{sec:mathspec}

В этом разделе из уравнений генераторов~\eqref{eq:a1}--\eqref{eq:a3} выводятся линейное и квадратичное ядра. Линейное ядро описывает окрестность покоя, а квадратичное задает первый нелинейный член разложения. Их проекции на дизайн эксперимента определяют, какая часть расхождения переходит в $\beta$ и контраст.

Обозначим шестимерное состояние $x=(z,f,v,q,x_E,x_I)^\top$, покой $x_0=(0,1,1,1,0,0)^\top$, систему $\dot x=F(x)+\mathbf b u$ и наблюдение $y=G(x)$, где $\mathbf b=(0,0,0,0,1,0)^\top$. Пусть $J=DF(x_0)$ --- гурвицева матрица Якоби (все ее собственные значения лежат в левой полуплоскости), $d=DG(x_0)$ --- строка производных наблюдения, $H_F=D^2F(x_0)$ и $H_G=D^2G(x_0)$ --- гессианы. При гладкости в окрестности покоя и нулевой предыстории точное первое ядро имеет вид
\begin{equation}\label{eq:h1}
h_1(\tau)=d\,e^{J\tau}\mathbf b,\quad \tau\ge0;\qquad h_1(\tau)=0,\quad\tau<0.
\end{equation}
Непрерывное симметричное ядро второго порядка определяется суммой
\begin{equation}\label{eq:k2}
k_2(\tau_1,\tau_2)=k_{2,\mathrm{field}}(\tau_1,\tau_2)+k_{2,\mathrm{obs}}(\tau_1,\tau_2),
\end{equation}
\begin{equation}\label{eq:k2field}
k_{2,\mathrm{field}}(\tau_1,\tau_2)=\frac12 d\int_0^{\min(\tau_1,\tau_2)}e^{Js}H_F[e^{J(\tau_1-s)}\mathbf b,\,e^{J(\tau_2-s)}\mathbf b]\,ds,
\end{equation}
\begin{equation}\label{eq:k2obs}
k_{2,\mathrm{obs}}(\tau_1,\tau_2)=\frac12[e^{J\tau_1}\mathbf b]^\top H_G[e^{J\tau_2}\mathbf b].
\end{equation}
При отрицательном хотя бы одном лаге $k_2=0$. Множитель $1/2$ уже включен в ядро, поэтому оператор
\begin{equation}\label{eq:Q}
Q[u_1,u_2](t)=\int_0^\infty\int_0^\infty k_2(\tau_1,\tau_2)\,u_1(t-\tau_1)\,u_2(t-\tau_2)\,d\tau_1d\tau_2
\end{equation}
не требует дополнительного внешнего множителя. Гессианы разделяют четыре источника квадратичной нелинейности: степенной закон оттока $v^{1/\alpha}$ в $\dot v$, нелинейное извлечение кислорода $g(f)=f[1-(1-E_0)^{1/f}]/E_0$ в $\dot q$, конвективный перенос дезоксигемоглобина $f_{\mathrm{out}}q/v$ и нелинейное уравнение наблюдения BOLD. Линейные E/I-блок и сопряжение с притоком не добавляют собственных квадратичных членов.

Статическая проекция определяется в обычном скалярном произведении $L^2$ на квадрате лагов $[0,T]^2$, $T=64$~с:
\begin{equation}\label{eq:kappa}
\kappa_0=\frac{\langle k_2,h_1\otimes h_1\rangle}{\|h_1\otimes h_1\|^2},\qquad
k_{2,\perp}=k_2-\kappa_0\, h_1\otimes h_1.
\end{equation}

Для разложения по событиям положим $u=\sum_{i\in A}u_i+\sum_{j\in B}u_j$, включив амплитуду и масштаб в каждое $u_i$. Квадратичная часть распадается на пять компонентов:
\begin{equation}\label{eq:five}
y_2=S_A+S_B+P_{AA}+P_{BB}+P_{AB},
\end{equation}
\begin{equation}\label{eq:fiveb}
\begin{split}
S_A&=\sum_{i\in A}Q[u_i,u_i],\qquad S_B=\sum_{i\in B}Q[u_i,u_i],\\
P_{AA}&=2\sum_{i<j,\;i,j\in A}Q[u_i,u_j],
\end{split}
\end{equation}
и аналогично $P_{BB}$, $P_{AB}=2\sum_{i\in A,j\in B}Q[u_i,u_j]$. Полный нелинейный остаток во временной области есть $r_{\mathrm{nl}}=y_2+r_{\ge3}$, где $r_{\ge3}$ объединяет вклады третьего и высших порядков. Проекция каждого слагаемого на тот же дизайн дает его вклад в ошибку контраста.

Для полной FV квадратичный оператор вычисляется по номинальной P-DCM на сетке лагов с шагом $0{,}05$~с в окне $64$~с; в подстановке используется полный квадратичный отклик, а не только статическая проекция.
